\chapter{Indigestion (Dyspepsia)}

\section*{Definition}
Indigestion (dyspepsia) is a group of upper abdominal symptoms including pain, burning, fullness, and nausea after eating. It affects approximately 20 percent of the population and is often functional (no identifiable organic cause).

\section*{What Happens in Your Body}
Dyspepsia results from a complex interaction of gastric acid, impaired gastric accommodation, delayed gastric emptying, visceral hypersensitivity, and psychological factors. Helicobacter pylori infection, NSAID use, and gastritis contribute in a subset of people. Functional dyspepsia involves abnormal brain-gut signaling without structural abnormality.

\section*{Causes}
\begin{itemize}
\item Functional dyspepsia (most common, no organic cause found)
\item Gastroesophageal reflux disease (GERD)
\item Peptic ulcer disease (gastric or duodenal ulcer)
\item H. pylori infection
\item Gastritis (acute or chronic)
\item Medications: NSAIDs, aspirin, corticosteroids, antibiotics
\item Gallbladder disease (biliary colic)
\item Pancreatitis
\item Gastric cancer (rare, but consider in older people with warning signs)
\item Delayed gastric emptying (gastroparesis)
\item Anxiety and stress
\item Food intolerances
\end{itemize}

\section*{When to See a Doctor}
\begin{itemize}
\item Upper abdominal pain or burning
\item Feeling of fullness after a small meal (early satiety)
\item Uncomfortable fullness after eating
\item Nausea, belching, bloating
\item Heartburn or regurgitation
\item Symptoms for 4 or more weeks
\item Unexplained weight loss (alarm symptom)
\item Vomiting, especially with blood
\item Difficulty swallowing
\item Iron deficiency anemia
\end{itemize}

\section*{Related Symptoms}
\begin{itemize}
\item Heartburn (GERD)
\item Bloating and gas
\item Nausea
\item Vomiting
\item Abdominal pain
\item Belching
\end{itemize}

\section*{Self-Care / Home Management}
\begin{itemize}
\item Eat smaller, more frequent meals
\item Avoid trigger foods (spicy, fatty, acidic)
\item Eat slowly and chew thoroughly
\item Avoid eating within 2-3 hours of bedtime
\item Elevate the head of the bed
\item Reduce or stop smoking
\item Limit alcohol and caffeine
\item Maintain a healthy weight
\item Manage stress through relaxation techniques
\end{itemize}

\section*{How Your Doctor Will Evaluate You}
\begin{itemize}
\item Clinical history and physical examination
\item Test for H. pylori (stool antigen or breath test)
\item Upper endoscopy (EGD) for people over 60 or with alarm symptoms
\item Abdominal ultrasound if biliary or pancreatic cause suspected
\item Gastric emptying study if gastroparesis suspected
\item Ambulatory pH monitoring if refractory GERD
\end{itemize}

\section*{Treatment Options}
\begin{itemize}
\item Antacids for mild symptoms
\item H2 receptor antagonists (famotidine)
\item Proton pump inhibitors (omeprazole, pantoprazole)
\item H. pylori eradication therapy if positive
\item Prokinetic agents (metoclopramide, domperidone) for gastroparesis
\item Tricyclic antidepressants (amitriptyline) for functional dyspepsia
\item Cognitive behavioral therapy for functional dyspepsia
\end{itemize}

\section*{References}

\begin{thebibliography}{99}
\bibitem{indigestion:romeiv-dyspepsia} Stanghellini V, et al. Rome IV: functional dyspepsia. \emph{Gastroenterology}. 2022;162(3):811-826.
\bibitem{indigestion:harrison-dyspepsia} Longo DL, et al. \emph{Harrison's Principles of Internal Medicine}. 21st ed. McGraw-Hill; 2022. Chapter 64: Dyspepsia.
\bibitem{indigestion:uptodate-dyspepsia} Talley NJ, et al. Functional dyspepsia in adults. In: UpToDate, Post TW (Ed), UpToDate, Waltham, MA; 2025.
\bibitem{indigestion:talley-lancet} Talley NJ, et al. Dyspepsia: a clinical review. \emph{Lancet}. 2023;401(10375):548-561.
\bibitem{indigestion:moayyedi} Moayyedi P, et al. Management of dyspepsia: a clinical practice guideline. \emph{BMJ}. 2024;384:e076507.
\bibitem{indigestion:suzuki} Suzuki H, et al. Pathophysiology of functional dyspepsia. \emph{J Gastroenterol}. 2023;58(4):321-334.
\end{thebibliography}
